\documentclass[12pt]{article}
\usepackage{amsmath}
\usepackage{amssymb}
\usepackage{physics}
\usepackage{graphicx}
\usepackage{hyperref}

\newcommand{\ds}{\displaystyle}

\title{Dynamics in Rotating Reference Frames}
\author{}
\date{}

\begin{document}
\maketitle

% ============================================================
\section{Introduction}
% ============================================================

This essay develops the dynamics of particles as observed in a rotating reference frame.
We begin by recalling the kinematical relations required
to compare time derivatives in inertial and rotating frames, then derive the transformation
 of the acceleration and the inertial forces that arise from it.

\subsection{Kinematical Preliminaries}

Let $\vb{A}(t)$ be the orthogonal matrix that maps the components of a vector in the inertial
frame to its components in the rotating frame:

\begin{equation}
  \vb{g}'(t) = \vb{A}(t)\,\vb{g}(t),
  \qquad
  \vb{g}(t) = \vb{A}^{\top}(t)\,\vb{g}'(t).
  \label{rotating-frames:frame-transformation}
\end{equation}

The angular velocity vector is defined as the time derivative of the infinitesimal rotation vector,

\begin{equation}
  \boldsymbol{\omega} = \lim_{\Delta t\to0}\frac{d\boldsymbol{\Omega}}{\Delta t}.
  \label{rotating-frames:angular-velocity-definition}
\end{equation}

The corresponding transformation law for time derivatives is

\begin{equation}
  \frac{d\vb{g}}{dt}
  = \vb{A}^{\top}(t)\!\left(
      \boldsymbol{\omega}\times\vb{g}' + \frac{d\vb{g}'}{dt}
    \right).
  \label{rotating-frames:time-derivative-rule}
\end{equation}


These equations summarize all the results from rigid-body kinematics required in the remainder of this essay.

% ============================================================
\section{Equations of Motion in a Rotating Frame}
% ============================================================

In this section we derive the equations of motion of a particle as observed in a rotating frame of reference
with respect to an inertial one.We restrict our discussion to the case of constant angular velocity. The resulting equations
 introduce the Coriolis and centrifugal forces that arise in a rotating frame.\par

We begin with the transformation law for time derivatives, Eq. \eqref{rotating-frames:time-derivative-rule}.
Differentiating once more gives the second derivative of the position vector in the inertial frame,

\[
  \frac{d^2\vb{r}(t)}{dt^2} =
    \frac{d\vb{A}^\top}{dt}\!\left(\boldsymbol{\omega}\times\vb{r}' + \frac{d\vb{r}'}{dt}\right)
    + \vb{A}^\top\!\left(\boldsymbol{\omega}\times\frac{d\vb{r}'}{dt} + \frac{d^2\vb{r}'}{dt^2}\right)
\]

From the orthogonality of the rotation matrix, $\vb{A}^\top\vb{A} = \vb{I}$:
\[
  \frac{d\vb{A}^\top}{dt}\,\vb{A} + \vb{A}\,\frac{d\vb{A}^\top}{dt} = 0
\]
we obtain:
\[
  \frac{d\vb{A}^\top}{dt} = -\vb{A}^\top\frac{d\vb{A}}{dt}\,\vb{A}^\top
\]

Inserting this expression in the equation above, and using the invariance of the cross product under proper orthogonal transformations,,

\[
  \begin{aligned}
    \frac{d^2\vb{r}(t)}{dt^2}
    &= -\vb{A}^\top\frac{d\vb{A}}{dt}\,\vb{A}^\top
       \!\left(\boldsymbol{\omega}\times\vb{r}' + \frac{d\vb{r}'}{dt}\right)
       + \vb{A}^\top\!\left(\boldsymbol{\omega}\times\frac{d\vb{r}'}{dt} + \frac{d^2\vb{r}'}{dt^2}\right) \\
    &= -\vb{A}^\top\frac{d\vb{A}}{dt}
       \!\left(\vb{A}^\top\boldsymbol{\omega}\times\vb{A}^\top\vb{r}'
         + \vb{A}^\top\frac{d\vb{r}'}{dt}\right)
       + \vb{A}^\top\!\left(\boldsymbol{\omega}\times\frac{d\vb{r}'}{dt} + \frac{d^2\vb{r}'}{dt^2}\right) \\
    &= -\vb{A}^\top\!\left(\vb{A}\!\left(\vb{A}^\top\boldsymbol{\omega}\times\vb{A}^\top\vb{r}'
         + \vb{A}^\top\frac{d\vb{r}'}{dt}\right)\right)\times\boldsymbol{\omega}
       + \vb{A}^\top\!\left(\boldsymbol{\omega}\times\frac{d\vb{r}'}{dt} + \frac{d^2\vb{r}'}{dt^2}\right) \\
    &= -\vb{A}^\top\!\left(\boldsymbol{\omega}\times\vb{r}' + \frac{d\vb{r}'}{dt}\right)\times\boldsymbol{\omega}
       + \vb{A}^\top\!\left(\boldsymbol{\omega}\times\frac{d\vb{r}'}{dt} + \frac{d^2\vb{r}'}{dt^2}\right)
  \end{aligned}
\]

we arrive at the equation:

\[
  \vb{A}\,\frac{d^2\vb{r}(t)}{dt^2} =
    \boldsymbol{\omega}\times(\boldsymbol{\omega}\times\vb{r}')
    + 2\boldsymbol{\omega}\times\frac{d\vb{r}'}{dt}
    + \frac{d^2\vb{r}'}{dt^2}.
\]

Introducing the inertial acceleration $\vb{a} = \frac{d^2\vb{r}(t)}{dt^2}$, and
the acceleration measured in the rotating frame as $\vb{a}' = \frac{d^2\vb{r}'(t)}{dt^2}$, we obtain
the acceleration transformation:

\begin{equation}
  \vb{a}' = -\boldsymbol{\omega}\times(\boldsymbol{\omega}\times\vb{r}')
    - 2\boldsymbol{\omega}\times\frac{d\vb{r}'}{dt}
    + \vb{A}\,\vb{a}
  \label{rotating-frames:acceleration-transformation}
\end{equation}

Equation \eqref{rotating-frames:acceleration-transformation} is purely kinematical.
No dynamical assumptions have been introduced yet. By applying Newton's second law
to the rotating frame, we identify

\[
\mathbf F' = m\mathbf a',
\]

and Eq. \eqref{rotating-frames:acceleration-transformation} becomes

\begin{equation}
  \vb{F}' = -m\boldsymbol{\omega}\times(\boldsymbol{\omega}\times\vb{r}')
    - 2 m \boldsymbol{\omega}\times\frac{d\vb{r}'}{dt}
    + m\vb{A}\,\vb{a}
  \label{rotating-frames:forces-transformation}
\end{equation}

In this equation there are terms that originate entirely from the kinematics of
the rotating frame. They are interpreted as apparent forces and are responsible
for the distinctive dynamics observed in non-inertial reference frames.

\subsection{The Centrifugal Force}

The centrifugal force is the first term of equation \eqref{rotating-frames:forces-transformation},

\begin{equation}
  \vb{F}_\text{cf} = -m\boldsymbol{\omega}\times(\boldsymbol{\omega}\times\vb{r}')
    = m\bigl( \omega^2\vb{r}'-\boldsymbol{\omega}\ip{\boldsymbol{\omega}}{\vb{r}'} \bigr)
  \label{rotating-frames:centrifugal-force}
\end{equation}

which has a magnitude proportional to the square of the angular velocity, and is directed perpendicular to
 the rotation axis and away from it. This force depends only on the angular velocity of the rotating frame
  and on the position of the body. It is independent of the body's velocity in the rotating frame. \par

The centrifugal force admits the potential
\[
  V_\text{cf}
  =
  \frac{m}{2}
  \left(
    -\ip{\vb r'}{\vb r'}\,
    \ip{\boldsymbol\omega}{\boldsymbol\omega}
    +
    \ip{\boldsymbol\omega}{\vb r'}^2
  \right),
\]

and it is therefore a conservative force, since

\[
  \vb F_\text{cf}
  =
  -\grad_{\vb r'} V_\text{cf}.
\]

\subsection{The Coriolis Force}

The Coriolis force is the term proportional to the velocity of the body in the
rotating frame appearing in equation \eqref{rotating-frames:forces-transformation},


\begin{equation}
  {\vb F}_\text{cor} = -2m\boldsymbol{\omega}\times\frac{d\vb{r}'}{dt}
  \label{rotating-frames:coriolis-force}
\end{equation}

This force arises only from the motion of the body with respect to the rotating
frame. Since it is always perpendicular to the velocity, it performs no work on
the body. \par


The Coriolis force has the same mathematical structure as the Lorentz force in
electromagnetism. Although it is not conservative, since it depends explicitly on the velocity,
it can nevertheless be derived from the generalized potential $V_\text{cor}$ determined as,

\[
  V_\text{cor} = -m\ip{\boldsymbol{\omega}\times\vb{r}'}{\dot{\vb{r}'}}
\]

so that the Coriolis force derives from a generalized potential:

\[
  \vb{F}_\text{cor} = \frac{d}{dt}\frac{\partial V_\text{cor}}{\partial\dot{\vb{r}'}}
    - \frac{\partial V_\text{cor}}{\partial\vb{r}'}
\]

To verify that the generalized potential indeed reproduces the Coriolis force,
we use the cyclic permutation identity of the scalar triple product,

\[
  \ip{\vb a\times\vb b}{\vb c}
  =
  \ip{\vb b\times\vb c}{\vb a}
  =
  \ip{\vb c\times\vb a}{\vb b},
\]

to rewrite the scalar product so that the vector with respect to which the
gradient is taken appears directly in the scalar product rather than in the
cross product. This allows the required derivatives with respect to
$\vb r$ and $\dot{\vb r}$ to be computed immediately. \par

Unlike an ordinary potential, a generalized potential may depend explicitly on
both the position and the velocity of the body. Such potentials naturally arise
in systems whose forces are linear in the velocity, such as the Coriolis force
and the Lorentz force in electromagnetism.

\subsection{Lagrange Equations of Motion in a Rotating Frame}

In the remainder of this section, we will denote the time derivative $da/dt$ to
$\dot{a}$ to simplify the equations. Unless otherwise stated, all vectors will be expressed in the rotating frame,
and $\vb r$ will denote the position vector in that frame. In this section
we would like to briefly touch on the functional form of the Lagrangian
 in the rotating frame of a body subject to \textbf{zero net force  in the inertial frame}  \par

The Lagrangian is obtained by adding to the kinetic energy all potential contributions, including generalized potentials when present.

\[
  L = T - V
\]

As the centrifugal and Coriolis forces admit potentials,  the Lagrangian of the system reads:

\[
  L = \frac{1}{2}m\ip{\dot{\vb{r}}}{\dot{\vb{r}}}
    +\frac{m}{2}\bigl(
      \ip{\vb{r}}{\vb{r}}\ip{\boldsymbol{\omega}}{\boldsymbol{\omega}}
      - \ip{\boldsymbol{\omega}}{\vb{r}}^2
    \bigr)
    + m\ip{\boldsymbol{\omega}\times\vb{r}}{\dot{\vb{r}}}.
\]

The Lagrange equations of motion are given by:

\[
  \frac{d}{dt}\frac{\partial L}{\partial\dot{\vb{r}}} - \frac{\partial L}{\partial\vb{r}} = 0
\]

The first term of the Lagrangian gives:

\[
  \frac{d}{dt}\frac{\partial L}{\partial\dot{\vb{r}}} = m\ddot{\vb{r}} + m{\boldsymbol \omega} \times \dot{\vb{r}}
\]

The second term of the Lagrangian gives:

\[
  \frac{\partial L}{\partial\vb{r}} = m (\boldsymbol{\omega}^2 {\vb r} -\ip{\boldsymbol{\omega}}{\vb r}\boldsymbol{\omega} ) +
  m\dot{\vb{r}}\times \boldsymbol{\omega}
\]
The Lagrange equations of motion in the rotating frame are therefore

\[
  m\ddot{\vb r}
  + 2m\boldsymbol{\omega}\times\dot{\vb r}
  + m\left(
    \ip{\boldsymbol{\omega}}{\vb r}\boldsymbol{\omega}
    - \omega^2\vb r
  \right)
  = 0.
\]

Dividing by $m$, we obtain

\[
  \ddot{\vb r}
  =
  -2\boldsymbol{\omega}\times\dot{\vb r}
  + \omega^2\vb r
  - \ip{\boldsymbol{\omega}}{\vb r}\boldsymbol{\omega},
\]

which is identical to equation
\eqref{rotating-frames:acceleration-transformation}.
\par

It has to be noted that a conservative force in the inertial frame is not necessarily
conservative in the rotating frame. A simple example can be easily constructed by considering
a constant force in the plane $z=0$ of the inertial frame such as:

\[
  \vb{F} = F_x\,\vb{i} + F_y\,\vb{j}
\]

where $F_x$ and $F_y$ are constants and $\vb{i}$, $\vb{j}$ are the unit vectors of the
inertial frame. The force in a frame rotating with constant angular velocity $\omega$ around
the axis $z$ is given by:

\[
  \vb{F}' =
  \begin{pmatrix}
    \cos\omega t  & \sin\omega t & 0 \\
    -\sin\omega t & \cos\omega t & 0 \\
    0             & 0            & 1
  \end{pmatrix}
  \vb{F}
  = (F_x\cos(\omega t) + F_y\sin(\omega t))\,\vb{i}'
  + (-F_x\sin(\omega t) + F_y\cos(\omega t))\,\vb{j}'
\]

where $\vb{i}'$, $\vb{j}'$ are the unit vectors of the rotating frame. The force in the
rotating frame is time-dependent, and does not have a generalized potential as it does not
explicitly depend on the velocity. Each case will have to be considered individually to
determine whether the dynamics can be described in a Lagrangian formalism, and in general the
equations of motion \eqref{rotating-frames:acceleration-transformation} will have to be used.

The existence of a Lagrangian for the inertial forces provides an elegant alternative formulation
of the dynamics in rotating reference frames. Nevertheless, when additional external forces are present,
each case must be analyzed individually to determine whether a Lagrangian description exists
or whether the equations of motion must instead be used directly.

% ============================================================\

\section{Life in a Rotating Reference Frame}

The dynamics of rotating reference frames is often counterintuitive, and even situations with which
we are very familiar become surprisingly subtle when examined in detail. One such example is the carousel. \par

If you are an observer at rest in the inertial frame fixed to the Earth, life is quite simple for you. All people standing in front of
the carousel are not moving and you see the people on the carousel going moving in circles around the centre of
rotation of the carousel, which has an angular velocity $\boldsymbol{\omega}$ parallel to the axis perpendicular to the Earth
surface, which we call ${\vb e}_z$. \par

If you are on the carousel, reality is more complicated
 First, you have to hold on to something in order not to fly away. You clearly perceive and measure a force that pushes
 you away from the center of rotation, whose magnitude increases with your distance from the center.

Furthermore, you see the people standing on the Earth moving around you in circles with constant angular velocity.
 This is quite a strange observation, because you cannot explain why they are not pushed away by the same force that you yourself must resist in order not to fly away.
  Even more interestingly, if something anchored to the carousel for
some reason slips away, it also starts to go in circle although following a more complicated dynamics as the people on the Earth.\par

The resolution of these apparent paradoxes lies in the combined action of the centrifugal and Coriolis forces in the rotating frame.\par

Consider a person standing on the Earth. Since the Earth is at rest in the
inertial frame, its position vector $\vb g(t)$ is constant in time, and therefore

\[
  \frac{d\vb g}{dt}=0.
\]

Using the transformation law relating the components of a vector in the fixed
and rotating frames,

\[
  \frac{d\vb g}{dt}
  =
  A^{-1}(t)
  \left(
    \boldsymbol{\omega}\times\vb g'
    +
    \frac{d\vb g'}{dt}
  \right),
\]

we conclude that

\[
  \frac{d\vb g'}{dt}
  =
  -\,\boldsymbol{\omega}\times\vb g'.
\]

Thus, an observer on the carousel sees every point fixed on the Earth moving
with a velocity equal to the opposite of the local rotational velocity of the
carousel.

The differential equation

\[
  \frac{d\vb g'}{dt}
  =
  -\boldsymbol{\omega}\times\vb g'
\]

can be solved explicitly. We seek a solution of the form

\[
  \vb g'(t)
  =
  \vb k\,e^{-i\alpha t},
\]

where $\vb k$ is a complex constant vector. Substituting into the differential
equation gives

\[
  i\alpha\vb k
  =
  -\boldsymbol{\omega}\times\vb k.
\]

Writing the complex vector as

\[
  \vb k
  =
  \vb a+i\vb b,
\]

with $\vb a$ and $\vb b$ real vectors, we obtain

\[
  i\alpha(\vb a+i\vb b)
  =
  -\boldsymbol{\omega}\times(\vb a+i\vb b),
\]

which, upon separating the real and imaginary parts, yields

\[
\begin{aligned}
  \alpha\vb a &= -\,\boldsymbol{\omega}\times\vb b,
  \\
  \alpha\vb b &= \phantom{-}\boldsymbol{\omega}\times\vb a.
\end{aligned}
\]

The second equation implies

\[
  \vb b
  =
  \frac{1}{\alpha}
  \boldsymbol{\omega}\times\vb a.
\]

 Substituting the expression for $\vb b$ into the first equation yields

\[
  \alpha\vb a
  =
  -\frac{1}{\alpha}
  \boldsymbol{\omega}\times
  \left(
    \boldsymbol{\omega}\times\vb a
  \right),
\]

or equivalently,

\[
  \alpha^2\vb a
  =
  -\boldsymbol{\omega}\times
  \left(
    \boldsymbol{\omega}\times\vb a
  \right).
\]

Using the vector identity

\[
  \vb u\times(\vb v\times\vb w)
  =
  \ip{\vb u}{\vb w}\vb v
  -
  \ip{\vb u}{\vb v}\vb w,
\]

we obtain

\[
  \alpha^2\vb a
  =
  -\left(
    \ip{\boldsymbol{\omega}}{\vb a}\boldsymbol{\omega}
    -
    \omega^2\vb a
  \right),
\]

which may be rewritten as

\[
  \left(
    \alpha^2-\omega^2
  \right)\vb a
  +
  \ip{\boldsymbol{\omega}}{\vb a}\boldsymbol{\omega}
  =
  0.
\]

Since the two vectors $\vb a$ and $\boldsymbol{\omega}$ are linearly
independent, the above equation implies

\[
  \alpha^2=\omega^2,
  \qquad
  \ip{\boldsymbol{\omega}}{\vb a}=0.
\]

Hence,

\[
  \alpha=\pm\omega,
\]

The general solution is therefore given by

\[
  \vb g'(t)
  =
  \left(
    \vb a
    +
    \frac{i}{\alpha}\,
    \boldsymbol{\omega}\times\vb a
  \right)
  e^{-i\alpha t}
  +
  \vb g_0',
  \qquad
  \alpha=\pm\omega,
\]

where $\vb g_0'$ is any constant vector parallel to
$\boldsymbol{\omega}$.

To obtain a more explicit expression, we define the orthonormal basis

\[
  \vb a=a\,\vb e_x,
  \qquad
  \boldsymbol{\omega}=\omega\,\vb e_z,
\]

where $\vb e_x$ is chosen along the direction of $\vb a$ and
$\vb e_z$ is parallel to the axis of rotation. Since

\[
  \vb e_y=\vb e_z\times\vb e_x,
\]

we obtain

\[
  \boldsymbol{\omega}\times\vb a
  =
  \omega a\,\vb e_y,
\]

and the solution becomes

\[
  \vb g'(t)
  =
  a
  \left(
    \vb e_x
    +
    i\frac{\omega}{\alpha}\vb e_y
  \right)
  e^{-i\alpha t}
  +
  \vb g_0'.
\]

The physical solutions are obtained by taking the real part of the complex
solutions. Using

\[
  e^{-i\alpha t}
  =
  \cos(\alpha t)-i\sin(\alpha t),
\]

we find

\[
\begin{aligned}
  \vb g'(t)
  &=
  \Re\left[
    a\left(
      \vb e_x
      -
      i\frac{\omega}{\alpha}\vb e_y
    \right)
    e^{-i\alpha t}
  \right]
  +
  \vb g'_0
  \\
  &=
  a\left[
    \cos(\alpha t)\vb e_x
    -
    \frac{\omega}{\alpha}
    \sin(\alpha t)\vb e_y
  \right]
  +
  \vb g'_0,
\end{aligned}
\]

where $\vb g'_0$ is constant and parallel to the axis of rotation.

For $\alpha=\omega$, this gives

\[
  \vb g'(t)
  =
  a\left[
    \cos(\omega t)\vb e_x
    -
    \sin(\omega t)\vb e_y
  \right]
  +
  \vb g'_0.
\]

For $\alpha=-\omega$, we obtain

\[
\begin{aligned}
  \vb g'(t)
  &=
  a\left[
    \cos(-\omega t)\vb e_x
    -
    \frac{\omega}{-\omega}
    \sin(-\omega t)\vb e_y
  \right]
  +
  \vb g'_0
  \\
  &=
  a\left[
    \cos(\omega t)\vb e_x
    -
    \sin(\omega t)\vb e_y
  \right]
  +
  \vb g'_0.
\end{aligned}
\]

Thus, both values of $\alpha$ produce the same real solution.

The same result can be obtained directly from the rotation matrix. Consider,
for simplicity, a point whose position in the fixed frame is

\[
  \vb g=a\vb e_x+\vb g_0,
\]

where $\vb g_0$ is parallel to the axis of rotation. Its components in the
rotating frame are obtained by applying the rotation matrix $A(t)$:

\[
\begin{aligned}
  \vb g'(t)
  &=
  A(t)\vb g
  \\
  &=
  a
  \begin{pmatrix}
    \cos(\omega t) & \sin(\omega t) & 0 \\
    -\sin(\omega t) & \cos(\omega t) & 0 \\
    0 & 0 & 1
  \end{pmatrix}
  \begin{pmatrix}
    1\\
    0\\
    0
  \end{pmatrix}
  +
  \vb g_0
  \\
  &=
  a\left[
    \cos(\omega t)\vb e_x
    -
    \sin(\omega t)\vb e_y
  \right]
  +
  \vb g_0.
\end{aligned}
\]

This is precisely the same solution obtained from the differential equation

\[
  \frac{d\vb g'}{dt}
  =
  -\boldsymbol{\omega}\times\vb g'.
\]

Now that we have determined the trajectory of a person standing on the Earth
as observed from the rotating frame, we can compute the centrifugal force
acting on this person,

\[
  \vb F_\text{cf}
  =
  m\omega^2\vb g'(t),
\]

where we have used the fact that the solution satisfies

\[
  \ip{\boldsymbol{\omega}}{\vb g'}=0.
\]

The Coriolis force is,

\[
\begin{aligned}
  \vb F_\text{cor}
  &=
  -2m\boldsymbol{\omega}\times
  \frac{d\vb g'}{dt}
  \\
  &=
  2m\boldsymbol{\omega}\times
  \left(
    \boldsymbol{\omega}\times\vb g'
  \right),
\end{aligned}
\]

where we have used

\[
  \frac{d\vb g'}{dt}
  =
  -\boldsymbol{\omega}\times\vb g'.
\]

Since $\vb g'$ is orthogonal to $\boldsymbol{\omega}$, we obtain

\[
\begin{aligned}
  \vb F_\text{cor}
  &=
  -2m\omega^2\vb g'(t).
\end{aligned}
\]

The total apparent force acting on the person in the rotating frame is therefore

\[
\begin{aligned}
  \vb F_\text{tot}
  &=
  \vb F_\text{cf}
  +
  \vb F_\text{cor}
  \\
  &=
  m\omega^2\vb g'(t)
  -
  2m\omega^2\vb g'(t)
  \\
  &=
  -m\omega^2\vb g'(t).
\end{aligned}
\]

This is precisely the centripetal force required to keep the person on the Earth
moving on the circular trajectory observed from the rotating frame.\par

\subsection{A Final Comment}

At this point it is worth pausing to discuss a few important conceptual points
that, although implicit throughout our derivation, play a fundamental role in
the dynamics of rotating reference frames. \par


So far, we have assumed that Newton's first principle of dynamics holds in the
rotating frame. An observer standing on the carousel could verify this
experimentally by measuring the force required to keep a mass $m$ at rest at
different positions on the carousel. The measured force would be precisely the
centrifugal force. Likewise, the observer could measure the force required to
keep a body moving with constant velocity. Repeating this experiment for
different velocities would allow the observer to determine the combined action
of the centrifugal and Coriolis forces. The fact that such experiments can be
performed and lead to well-defined force fields is a remarkable property of
rotating frames and is by no means obvious. \par

Once these force fields have been determined - assuming no other forces are
present - the observer can predict the motion of any body from its initial
conditions by solving Newton's second law or, equivalently, the Lagrange
equations. Thus, in the rotating frame, Newton's second law, Hamilton's
principle, the Lagrange equations, and the Hamilton equations remain perfectly
valid descriptions of the dynamics. \par

What would nevertheless surprise the observer is that the apparent forces do
not satisfy Newton's third principle of dynamics. In particular, there is no
body in the universe on which acts a force equal and opposite to the
centrifugal force. This lack of an action-reaction pair is one of the reasons
why the centrifugal and Coriolis forces are referred to as \emph{apparent}
forces. \par

In classical mechanics, inertial reference frames are precisely those in
which Newton's three principles of dynamics hold. The rotating frame considered in this
chapter is therefore not an inertial frame. Remarkably, by introducing the centrifugal
and Coriolis forces, Newton's first and second principles can be restored, allowing
the dynamics to be described exactly as in an inertial frame. Newton's third principle,
however, cannot be recovered: there is no body on which acts a force equal and opposite to the centrifugal or Coriolis
force. This is one of the defining characteristics of apparent forces. \par

  One should not conclude that every force that fails to satisfy
  the simple action-reaction principle is therefore an apparent force.
   The Lorentz force on a charge moving in a magnetic field provides an important counterexample.
   Unlike the centrifugal and Coriolis forces, the Lorentz
  force cannot be eliminated by changing the reference frame. It represents a
   genuine physical interaction and therefore belongs to a completely different category. \par

There is, however, an even deeper reason why these forces are regarded as
apparent. The Lorentz force
 is fundamentally different from the centrifugal and Coriolis forces. No
change of coordinates can eliminate the Lorentz force, since it originates from
the electromagnetic interaction and, in special relativity, arises from the
transformation of the electric field of the charge in its rest frame. The
centrifugal and Coriolis forces, on the other hand, can be made to disappear
completely simply by observing the same motion from an inertial frame, such as
the frame of an observer standing on the Earth. It is this dependence on the
choice of reference frame that ultimately distinguishes apparent forces from
genuine physical interactions.


\end{document}
